% =====================================================================
\chapter{Introduction}
\label{ch:intro}

\section{Background}
For a large share of the students of any university, the day begins and
ends with a bus. University transport is often the only affordable way to
travel between home and campus, and in Bangladesh---where many students
commute considerable distances into district towns such as Manikganj---the
university bus is a daily necessity rather than a convenience. Yet the
information available to the rider about that bus has barely changed in
decades. The timetable on a notice board says when the bus is
\emph{supposed} to leave; nothing tells a student where the bus
\emph{actually} is right now.

Meanwhile, real-time vehicle tracking itself is a mature technology.
Satellite positioning receivers are inexpensive, cellular data coverage
reaches nearly every road a university bus travels, and every student
carries a smartphone capable of rendering a live map. Ride-sharing
platforms have made the ``watch your vehicle approach on the map''
experience an everyday expectation. The technical ingredients for live bus
tracking are therefore all available; what has been missing at
institutions like ours is a system that assembles those ingredients at a
cost a university transport office can justify, in a form students will
actually use, and in the language they speak.

\section{Motivation}
The motivation for this work came directly from our own daily experience
and that of our fellow students. A student who arrives at a roadside stop
faces an information vacuum with real costs:

\begin{itemize}
  \item \textbf{Missed buses.} If the bus passed early, the student may
  wait for a vehicle that has already gone, and discover it only when it
  is too late to arrange an alternative.
  \item \textbf{Wasted waiting time.} If the bus is late, the student
  stands at the stop---often without shelter, in heat or rain---for time
  that could have been spent at home or in the library had the delay been
  visible in advance.
  \item \textbf{Crowding and uncertainty at the origin.} Without knowing
  when a bus will actually depart, students cluster early around the
  boarding point, producing avoidable congestion.
\end{itemize}

Commercial fleet-tracking products could in principle address this, but
three barriers make them a poor fit for a university in a developing
region. First, their subscription pricing is set for commercial logistics
fleets, not for a cost-sensitive academic transport office. Second, they
are built for the fleet \emph{operator}: their web dashboards answer
``where are my vehicles?'' for a manager, not ``when will my bus reach my
stop?'' for a passenger. Third, their interfaces are almost always
English-only, whereas a large portion of our student population is more
comfortable in Bangla.

We therefore set out to build, deploy, and evaluate a complete system of
our own---not a simulation or a prototype confined to a laboratory, but a
platform running on the real university fleet with real riders.

\section{Problem Statement}
The problem addressed by this thesis can be stated as follows:

\begin{quote}
\emph{Students using university transport cannot observe the live position
of their bus or a dependable estimate of its arrival time at their own
stop. Existing commercial tracking solutions are too costly for
institutional adoption, are designed for fleet operators rather than
passengers, and are not localized. Furthermore, tracking approaches that
depend on a driver's smartphone are fragile---they fail whenever the
driver forgets, declines, or is unable to run the tracking application.}
\end{quote}

Decomposing this, the specific technical problems to solve were:

\begin{enumerate}
  \item \textbf{Autonomous position acquisition.} Obtain a continuous
  stream of bus positions without depending on any daily human action,
  while still tolerating the failure or absence of the tracking hardware.
  \item \textbf{Noise robustness.} Convert raw, noisy GPS fixes---which
  contain urban scatter, occasional impossible jumps, and stationary
  jitter---into a stable position that can be shown to the public.
  \item \textbf{Passenger-meaningful information.} Translate a raw
  position into what a rider actually needs: the arrival time at
  \emph{their} stop, the state of the bus \emph{before} departure, and an
  alert when their drop-off approaches.
  \item \textbf{Low-latency delivery at low cost.} Push every update to
  every interested phone within seconds, using infrastructure cheap enough
  for a student-project budget.
  \item \textbf{Localization.} Present the entire experience in both
  English and Bangla.
\end{enumerate}

\section{Objectives}
The objectives of this thesis work were set at the start of the project
and are listed here in the order in which the system must satisfy them:

\begin{enumerate}
  \item To design a three-tier, event-driven architecture for real-time
  university bus tracking that runs entirely on a single low-cost server.
  \item To implement resilient dual-source GPS acquisition, in which a
  fixed vehicle-grade GPS/GSM tracker acts as the fully automatic primary
  source and the driver's smartphone acts only as an optional automatic
  fallback, with per-fix health-based arbitration between them.
  \item To implement a location-quality filtering pipeline that gates
  every fix on accuracy, physical plausibility, and minimum movement
  before it can move the publicly visible position.
  \item To design a lightweight, training-data-free arrival-time estimator
  that produces a per-stop ETA with a confidence label on every accepted
  fix.
  \item To build a bilingual (English/Bangla) passenger mobile application
  that remains informative in edge conditions---before departure, during
  signal loss, and near the rider's stop---and a web administration
  console for routes, schedules, fleet, and live operations.
  \item To deploy the complete system on the operational university fleet
  and evaluate it over a multi-week pilot using real operational data.
\end{enumerate}

\section{Scope of the Work}
The scope of this thesis covers the complete life of a position report:
from the GPS receiver on the bus, through cellular transmission, decoding,
arbitration, filtering, ETA computation, storage, and publication, to the
pixel on a passenger's screen. It includes the hardware selection and
configuration of the on-vehicle tracker, the backend server and database,
the realtime delivery layer, the passenger application, the administration
console, and the evaluation of the deployed whole.

Certain related problems are deliberately out of scope. The system does
not attempt learning-based arrival prediction (it logs the data needed to
train such models later); it does not process payments or ticketing; it
does not measure passenger occupancy at scale; and the pilot was confined
to one university's fleet in and around Manikganj rather than a
multi-institution deployment. These boundaries are revisited as future
work in Chapter~\ref{ch:conclusion}.

\section{Contributions}
The principal contributions of this thesis are:

\begin{enumerate}
  \item \textbf{A deployed end-to-end architecture} for real-time
  university bus tracking centred on an autonomous fixed GPS tracker with
  an optional, health-arbitrated smartphone fallback, providing resilient
  coverage at low cost (Chapters~\ref{ch:design}
  and~\ref{ch:implementation}).
  \item \textbf{A lightweight geometric ETA estimator} that projects the
  live position onto the route polyline and derives per-stop arrival times
  from a rolling-average speed, annotated with a confidence level, and
  requiring no training data (Section~\ref{sec:design-eta}).
  \item \textbf{A passenger-facing design contribution}: a bilingual user
  experience with pre-trip vehicle awareness, staleness signalling, and
  proximity-based drop-off alerts, which keeps the interface informative
  precisely in the edge conditions where naive trackers show a blank map
  (Section~\ref{sec:impl-passenger}).
  \item \textbf{An operational evaluation} of the deployed system over a
  seven-week pilot on a real fleet, reporting acquisition latency,
  source-fallback usage, and passenger engagement
  (Chapter~\ref{ch:evaluation}).
  \item \textbf{The identification and elimination of idle-reporting
  waste}: pilot data showed that roughly 98\% of raw device reports were
  transmitted while buses were parked; a two-sided optimization (device
  reporting cadence and server polling cadence) cut idle-period reporting
  by up to $60\times$ without affecting live tracking
  (Section~\ref{sec:eval-efficiency}).
\end{enumerate}

\section{Organization of the Thesis}
The remainder of this thesis is organized as follows.
Chapter~\ref{ch:literature} reviews related work on campus bus tracking,
smartphone-based positioning, arrival-time prediction, GPS quality
handling, and open telematics platforms, and positions this work against
it. Chapter~\ref{ch:analysis} analyses the system's stakeholders and
requirements and records the feasibility study and development
methodology. Chapter~\ref{ch:design} presents the system architecture and
the design of each major component. Chapter~\ref{ch:implementation}
describes the implementation of the backend, the realtime layer, the two
client applications, and the on-vehicle hardware configuration.
Chapter~\ref{ch:evaluation} reports the seven-week pilot evaluation and
discusses its findings, including the idle-reporting optimization.
Chapter~\ref{ch:conclusion} concludes and outlines future work. The
appendices record the tracker's SMS configuration commands, the
operator-configurable parameters of the processing pipeline, and
representative screens of the passenger application.
